\par\bigskip
\SourceLeaf{489}%
\pdfbookmark[0]{Entwurf zur Gradmessung. 1819.}{entwurf-gradmessung-1819}
\begin{center}\large [ENTWURF ZUR GRADMESSUNG.]\end{center}
\par
— — —

Durch einen Brief von Professor \textsc{Schumacher} habe ich erfahren, dass der dänische Gesandte in London, Graf \textsc{Bourke}, während des Aufenthalts des Professors dort mit Graf \textsc{Münster} über den Nutzen der Fortsetzung der dänischen Gradmessung durch das Königreich Hannover gesprochen hat. Letzterer hat darum gebeten, dass ich ihm über diesen Gegenstand schreibe. Ich musste diese Mitteilung als einen Befehl auffassen. Das warme Interesse, das Ew. — — an dieser wissenschaftlichen Angelegenheit genommen und schon im vergangenen Jahr durch Maßnahmen, die die Möglichkeit einer solchen Fortsetzung sicherten, bekundet haben, verpflichtet mich, Ew. — — sofort davon zu unterrichten. Es gibt mir zugleich die Gewissheit, dass Sie meine Bitte, die Sache in London nach Möglichkeit zu unterstützen, wohlwollend aufnehmen werden. Ich darf umso weniger zögern, diese Bitte auszusprechen, als ich dabei einzig aus reiner Liebe zur Wissenschaft handle. Ohne diese Liebe und ohne die lebhafte Überzeugung von der großen Bedeutung eines solchen Unternehmens würde ich nach meiner persönlichen Neigung ein längeres Fernbleiben von den rein astronomischen Arbeiten eher als ein Opfer ansehen. Auch der ehrwürdige Sir \textsc{Joseph Banks} hat, wie mir Professor \textsc{Schumacher} schreibt, die Idee einer weiteren Fortsetzung dieser Gradmessung auf dem Kontinent mit großer Wärme aufgegriffen und jede Mitwirkung versprochen.

Professor \textsc{Schumacher} wird Ew. — — bei seiner Durchreise durch Hannover selbst gesagt haben, dass er nach England gereist ist, um den berühmten großen Zenitsektor in Empfang zu nehmen, der bei der englischen Gradmessung verwendet wurde und ihm für seine astronomischen Beobachtungen geliehen wird. Gleich nach seiner Rückkehr wird er am südlichen Endpunkt seiner Messungen in Lauenburg mit diesen Beobachtungen beginnen. Bei der Aussicht auf eine Fortsetzung der dänischen Gradmessung durch das Königreich Hannover wäre es zweifellos in vieler Hinsicht für beide Gradmessungen sehr wünschenswert, wenn ich \textit{diese} Beobachtungen gemeinsam mit ihm anstellte. \textit{Ohne} diese Zusammenarbeit ist es leicht denkbar, dass künftig Abweichungen oder Zweifel auftreten könnten, die sich dann nur schwer klären ließen und denen \textit{durch} diese Zusammenarbeit vorgebeugt würde. So sind ja gerade über die astronomischen Beobachtungen, die am nördlichen Endpunkt der französischen Gradmessung in Dünkirchen angestellt wurden, später Bedenken aufgekommen. Deshalb hielt man es für nötig, dass englische und französische Kommissare im vergangenen Herbst mit dem genannten Zenitsektor gemeinsam neue Beobachtungen anstellten. Und wie ich höre, soll eben dieser Zenitsektor, sobald Professor \textsc{Schumacher} ihn zurückgegeben hat, zu einer nochmaligen Wiederholung der Beobachtungen auf den Balearen dorthin verschifft werden. Nur wenn ich selbst mit diesem Instrument beobachte und diese Erfahrungen mit denen vergleiche, die ich am Meridiankreis von \textsc{Repsold} gemacht habe, und noch mehr mit denen, die ich künftig mit dem Kreis von \textsc{Reichenbach} machen werde, könnte ich beurteilen, ob eben dieser Zenitsektor für die künftigen Beobachtungen im Königreich Hannover notwendig oder besonders wünschenswert sein wird. Selbst wenn die Fortsetzung der Gradmessung nicht zustande käme, wären die Erfahrungen, die ich mit dem Zenitsektor, einem in seiner Art einzigartigen Instrument, machte, für meine künftigen Beobachtungen an den Meridianinstrumenten der hiesigen Sternwarte nicht ohne wesentlichen Nutzen. Als Beleg dafür, für wie wichtig man die Beobachtungen in Lauenburg hält, darf ich noch anführen, dass mir Doktor \textsc{Olbers} von seiner Absicht geschrieben hat, nach Lauenburg zu kommen, um das Instrument kennenzulernen und den Beobachtungen beizuwohnen.
